\newcommand*{\version}{3.33}

\documentclass[12pt,
              article,
              oneside
              ]{memoir}

\usepackage{iftex}

\ifluatex
  \usepackage{fontspec}
\fi

\usepackage[
  % latin,
  english,
  % spanish,% <- principal
  % es-minimal,
  % es-noshorthands,
  % es-tabla,
]{babel}

\ifluatex
% LuaLaTeX --------------------
\defaultfontfeatures{Ligatures=TeX, Scale=MatchLowercase}
\usepackage[MainFeatures={medium,
  Ligatures=Common,
  Numbers=OldStyle,
  SmallCapsFeatures = {LetterSpace=5.0}}]
{junicodevf}
%
% \setsansfont{Libertinus Sans}
% (without the ligatures of TeX, so that -- stays two hyphens)
\setmonofont[Scale=MatchLowercase, Ligatures=TeXOff]{DejaVu Sans Mono}
% \setmonofont[Scale=MatchLowercase]{Libertinus Mono}
% \setmonofont[Scale=MatchLowercase}]{Fira Code}
\else
% PDFLaTeX ---------------------
\usepackage[T1]{fontenc}
\usepackage[osf]{libertinus}
\usepackage[scaled=0.9]{FiraMono}
% \usepackage[scaled=0.9]{cascadia-code}
% \usepackage[scaled=0.9]{roboto}
\input glyphtounicode
\pdfgentounicode=1
\fi

\usepackage[babel=true, verbose=false, tracking=true, expansion=true, protrusion=true, final, draft=false]{microtype}
\SetTracking{ encoding = *, shape = sc }{ 30 }
\SetTracking[context = notracking, ]{encoding = *}{0}

\usepackage[autostyle=false, style=british]{csquotes}
\usepackage{latexcolors}

\newcommand{\Rojo}{\color[HTML]{8B0000}}
\newcommand*{\rojo}[1]{\textcolor[HTML]{8B0000}{#1}}
\newcommand*{\rojoit}[1]{\textit{\textcolor[HTML]{8B0000}{#1}}}
\newcommand{\Azul}{\color{bluenivaca}}
\newcommand*{\azul}[1]{\textcolor{bluenivaca}{#1}}
\newcommand{\Verde}{\color{darkgreen}}
\newcommand*{\verde}[1]{\textcolor{darkgreen}{#1}}
\newcommand{\Amarillo}{\color{goldenpoppy}}
\newcommand*{\amarillo}[1]{\textcolor{goldenpoppy}{#1}}
\newcommand{\Morado}{\color{byzantium}}
\newcommand*{\morado}[1]{\textcolor{byzantium}{#1}}
\newcommand{\Naranja}{\color{orange}}
\newcommand*{\naranja}[1]{\textcolor{orange}{#1}}

% \usepackage[margin=3cm]{geometry}

\usepackage{graphicx}
\usepackage{booktabs}

\clubpenalty=9999
\widowpenalty=9999
\brokenpenalty=5000
\frenchspacing
\raggedbottom
\midsloppy
\parindent=14pt
\fboxsep=2pt

\renewcommand\thesection{\arabic{section}}
\chapterstyle{reparticle}

\usepackage{enumitem}
\setlist{noitemsep, nosep}
\setlist[description]{font=\itshape}

\newcommand*{\elipsis}{\,.\,.\,.\,}
\newcommand*{\elipsisb}{\textup{[\elipsis]}}

\renewenvironment{quote}%
 {\begin{list}{}%
              {\setlength\rightmargin{0pt}%
              \setlength\leftmargin{\parindent}}%
  \item[]\small\ignorespaces}
 {\unskip\end{list}}


\definecolor{Link}{HTML}{408080}
\usepackage{hyperref}
\hypersetup{
  colorlinks=true,
  linkcolor=blue,
  filecolor=blue,
  citecolor = black,
  urlcolor=Link,
}
\urlstyle{tt}
\usepackage{bookmark}

\setlength{\droptitle}{-2cm}

\title{\bfseries%
  \texttt{py-words} v.\version{} \\
  User Manual}

\author{%
  Nicolás Vaughan\\
  \begin{minipage}[t]{0.8\linewidth}
    \centering\small%
    \href{mailto:nicolasvaughan@proton.me}{\texttt{nicolasvaughan@proton.me}}\\
    \url{https://nicolasvaughan.org}\\
    \url{https://github.com/nivaca/py-words}\\
    \url{https://codeberg.org/nivaca/py-words}\\
  \end{minipage}}


\date{2026-11-27}

% The pictures, made by tools/screenshots.py
\graphicspath{{img/}}
\usepackage{longtable}

% the title and the author in the PDF's own metadata
\hypersetup{pdftitle={py-words User Manual}, pdfauthor={Nicolás Vaughan}}

% the index (makeindex, with the style manual.ist: see latexmkrc)
\makeindex

% menus, keys, and the pictures of the window
\newcommand*{\menu}[1]{\textit{#1}}
\newcommand*{\then}{\,>\,}
\newcommand*{\key}[1]{\texttt{#1}}
\newcommand*{\shot}[3]{%
  \begin{figure}[htbp]
    \centering
    \includegraphics[width=#2\textwidth,height=0.75\textheight,keepaspectratio]{#1}
    \caption{#3}
    \label{fig:#1}
  \end{figure}}
% the tilde of the interactive commands (~E, ~L ...)
\newcommand*{\tl}{\textasciitilde}

\begin{document}
\maketitle

\tableofcontents*

\section{Introduction}
\noindent
\texttt{py-words} is a Latin dictionary that knows Latin grammar. It
analyses an inflected Latin word, or a whole line of Latin, into the
dictionary entries and the forms it can be, with their meanings in English
or in Spanish, and it finds the Latin for an English or a Spanish word. It
can be used in four ways: in a window, on the command line, in a terminal in
full screen, and in a web browser. It is also online, with nothing to
install, at \url{https://nivaca.pythonanywhere.com/}.

\index{WORDS (Whitaker)}\index{Lewis and Short}\index{De Miguel}\index{Echave-Sustaeta}\index{Latham}\index{Collatinus}%
I wrote it in 2026. Its starting point was WORDS, the program of William
Whitaker, but the code and the data were redesigned, and the dictionary
grew: it now has the rare words and the English definitions of Lewis and
Short's \textit{A Latin Dictionary}, the words, names and Spanish
definitions of R.~de~Miguel's \textit{Nuevo diccionario latino-español
etimológico}, the basic vocabulary of J.~Echave-Sustaeta, R.~E. Latham's
medieval word-list, and the long vowels of Collatinus. Each entry says
where it comes from (section~\ref{sec:codes}).

\index{conventions}\index{macos and windows@macOS and Windows}%
This manual explains first how to install and start the program, then the
window, how to search and how to read the results, and after that the
other three programs, the settings and the dictionary editors. The menus
are written in italics, with their items after a \enquote{>}, as
\menu{Dictionary\then Settings\ldots}; keys are written as \key{Ctrl+F}.
On macOS, \key{Ctrl} is the Command key. The pictures show the programs in
English; with the interface in Spanish, every menu, message and help text
is in Spanish too (section~\ref{sec:languages}).

\index{licence}%
\texttt{py-words} is free software, under the GNU General Public License,
version 3 or later. It uses only Python and its standard library (with Tk
for the windows and curses for the terminal), so it needs nothing else; the
full list of the software it uses, with its licences, is in
\texttt{docs/LICENCES.md}. WORDS and its data were given by Whitaker
\enquote{for any and all use}, and the text of Lewis and Short comes from
the Perseus Digital Library, under CC BY-SA 4.0.

\section{Installing it}
\label{sec:install}

Every version is published on GitHub and on Codeberg (the addresses are on
the first page), with these files:

\index{appimage@AppImage}\index{fedora (copr, rpm)@Fedora (Copr, RPM)}\index{debian and ubuntu (deb)@Debian and Ubuntu (.deb)}\index{zip}%
\begin{itemize}
\item \texttt{py-words-v.VERSION-windows-x64-setup.exe},
  \texttt{py-words-v.VERSION-macos-arm64.pkg} (Apple silicon),
  \texttt{py-words-v.VERSION-macos-x86\_64.pkg} (Intel) and
  \texttt{py-words-v.VERSION-linux-x86\_64.run}: the installers, for
  Windows, macOS and Linux;
\item \texttt{py-words-v.VERSION-x86\_64.AppImage}: the AppImage, for any
  Linux;
\item \texttt{py-words-VERSION-1.fc45.noarch.rpm}: the Fedora package,
  signed, with the key that signs it, \texttt{RPM-GPG-KEY-py-words};
\item \texttt{py-words\_VERSION-1\_all.deb}: the package for Debian and
  Ubuntu, with its signature, \texttt{py-words\_VERSION-1\_all.deb.asc};
\item \texttt{py-words-v.VERSION.zip}: the program itself, for Linux,
  macOS and Windows.
\item \texttt{SHA256SUMS}, with \texttt{SHA256SUMS.asc}: the SHA-256 sums of
  the files, and their signature. Every file also has its own signature
  beside it (\texttt{FILE.asc}), made with the same key, as the tag of the
  version is; once the key is imported (\texttt{gpg --import
  RPM-GPG-KEY-py-words}), \texttt{gpg --verify FILE.asc FILE} checks a file.
\end{itemize}

\subsection{What is needed}

\index{python@Python}\index{tk@Tk}%
Only Python 3.13 or later. The windows (the dictionary and its editor)
need Tk, which comes with Python on Windows and with the installers of
python.org on macOS. On Linux, and with Homebrew's Python on macOS, it may
be a separate package:

\begin{verbatim}
$ sudo dnf install python3-tkinter      (Fedora)
$ sudo apt install python3-tk           (Debian, Ubuntu)
$ brew install python-tk                (macOS, Homebrew)
\end{verbatim}

\noindent
\texttt{python3 --version} says which Python you have; if it is too old,
the program says so and stops. The commands in this manual use
\texttt{python3}; on Windows, write \texttt{python} instead. The
installers, the AppImage and the packages bring or use their own Python,
so they need nothing of this.

\subsection{The installers (Windows, macOS, Linux)}
\label{sec:installers}

\index{installers}\index{macos and windows@macOS and Windows}%
The simplest way. Each installer brings its own Python 3.13, with Tk, and
installs the program, the entries of the menu, the \texttt{py-words}
command (with the same first argument as the AppImage, below) and a way
to remove it all. A newer version replaces the older one; the settings,
the history and the cache are kept. As in the packages, the installed
dictionary cannot be edited (section~\ref{sec:editing} says how to do it
anyway).

\begin{description}
\item[Windows] (10 or later, 64-bit). Run the \texttt{.exe}. It installs
  for you alone, without administrator rights, unless you choose on its
  first page to install for all users. It speaks English or Spanish. It
  adds the window to the Start menu and, if you want, the command line,
  the full-screen terminal version and the web page, a shortcut on the
  desktop and the \texttt{py-words} command in the PATH. It is removed
  from ``Installed apps'' in the Settings, and asks then whether to remove
  your settings too. It is not signed, so SmartScreen may warn about it:
  ``More info'', then ``Run anyway''.
\item[macOS] (11 or later). Open the \texttt{.pkg} for your Mac (arm64
  for Apple silicon, x86\_64 for Intel). It puts \texttt{py-words.app} in
  Applications and, unless you leave them out under ``Customize'',
  \texttt{py-words (CLI)}, \texttt{(TUI)} and \texttt{(Web)}, which open in
  Terminal, and the \texttt{py-words} command. It is not signed by Apple,
  so macOS refuses it the first time: in System Settings, ``Privacy \&
  Security'', choose ``Open Anyway''. \texttt{sudo py-words uninstall}
  removes it (\texttt{--purge}: also your settings, history and cache).
\item[Linux] (x86-64). Run \texttt{sh py-words-v.VERSION-linux-x86\_64.run}.
  As a normal user it installs for you alone, in
  \texttt{\textasciitilde/.local/share/py-words}, with the command in
  \texttt{\textasciitilde/.local/bin}; with \texttt{sudo}, for everyone, in
  \texttt{/opt/py-words} and \texttt{/usr/local/bin}. It adds the four
  menu entries and the icons. It says first what it will do and asks:
  \texttt{--yes} asks nothing, \texttt{--only-gui} adds only the entry of
  the window, \texttt{--help} lists the other options.
  \texttt{py-words uninstall} removes it.
\end{description}

\subsection{The ZIP, or a copy of the repository}

\index{zip}\index{git}%
Unpack the ZIP anywhere, or clone the repository:

\begin{verbatim}
$ git clone https://github.com/nivaca/py-words.git
\end{verbatim}

\noindent
The folder is the whole program: there is nothing to build or to install,
and it runs from where it is.

\begin{verbatim}
$ cd py-words
$ python3 words.py amicus          (a first test)
$ python3 words_gui.py             (the window)
\end{verbatim}

\subsection{The AppImage (Linux)}

\index{appimage@AppImage}%
The AppImage is a single file that brings its own Python, with Tk. Make it
executable and run it:

\begin{verbatim}
$ chmod +x py-words-v.3.33-x86_64.AppImage
$ ./py-words-v.3.33-x86_64.AppImage              (the window)
$ ./py-words-v.3.33-x86_64.AppImage cli amicus   (the command line)
$ ./py-words-v.3.33-x86_64.AppImage help         (the others)
\end{verbatim}

\noindent\index{dictionary editors!read-only data}%
The first argument chooses the program: \texttt{gui} (the default),
\texttt{cli}, \texttt{tui}, \texttt{web}, \texttt{edit},
\texttt{edit-gui}, \texttt{check} or \texttt{convert}; what follows goes to
that program. The AppImage cannot be written to, so the dictionary cannot
be edited in it (section~\ref{sec:editing} says how to do it anyway).
Tools such as AppImageLauncher or Gear Lever add it to the menu of the
system.

\subsection{Fedora}

\index{fedora (copr, rpm)@Fedora (Copr, RPM)}%
The easiest way is the Copr repository. The package uses Fedora's own
Python and Tk, and adds four entries to the menu of the system, with the
icon:

\begin{verbatim}
$ sudo dnf copr enable nivaca/py-words
$ sudo dnf install py-words
\end{verbatim}

\noindent
Updates then come with \texttt{sudo dnf upgrade}. The RPM of each release
can also be installed by hand, after importing its key once with
\texttt{sudo rpm --import RPM-GPG-KEY-py-words}. It runs on any Fedora
with Python 3.13 or later, whatever the \texttt{fc} number in its name.
The command \texttt{py-words} then starts the programs, with the same
first argument as the AppImage:

\begin{verbatim}
$ py-words                 (the window)
$ py-words cli amicus      (the command line)
$ py-words help            (the other programs)
\end{verbatim}

\noindent\index{dictionary editors!read-only data}%
The program is installed in \texttt{/usr/share/py-words}, which cannot be
written to, so the dictionary cannot be edited there either, as in the
AppImage.

\subsection{Debian and Ubuntu}

\index{debian and ubuntu (deb)@Debian and Ubuntu (.deb)}%
The \texttt{.deb} package installs the same as the RPM, and also uses the
system's Python and Tk. It needs Python 3.13 or later (Debian 13, Ubuntu
25.04 or later); \texttt{apt} installs \texttt{python3-tk} with it:

\begin{verbatim}
$ sudo apt install ./py-words_3.33-1_all.deb
\end{verbatim}

\noindent
Its signature comes apart, in the \texttt{.deb.asc} file, made with the
same key as the RPM; \texttt{gpg --verify} checks it once the key is
imported. The command \texttt{py-words} then starts the programs, as with
the RPM, and \texttt{sudo apt remove py-words} removes it. There is no
repository: each new version is installed by hand in the same way.

\subsection{A menu entry}

\index{menu entries}\index{macos and windows@macOS and Windows}%
With the ZIP or a copy of the repository, run once

\begin{verbatim}
$ python3 words_install.py
\end{verbatim}

\noindent
to add the program to the menu of the system, with four entries: the
window, the command line and the full-screen terminal version (these two
open in a terminal), and the web version (its server runs in a terminal,
and the browser opens). On Linux they go to the applications menu, on
Windows to the Start menu, and on macOS to \texttt{\textasciitilde/Applications}.
\texttt{--only-gui} adds only the window, and \texttt{--uninstall} removes
them. It needs no administrator rights. The entries start the program in
its folder, so if the folder is moved, run it again.

\subsection{Updating and uninstalling}

\index{updating}\index{uninstalling}%
With git, \texttt{git pull} in the folder; if not, download the ZIP again
and replace the folder. If you edited the dictionary, keep a copy of
\texttt{data/dictionary.xml} first. To uninstall, remove the menu entries
with \texttt{python3 words\_install.py --uninstall} and delete the folder;
the settings, the history and the cache can be deleted too
(section~\ref{sec:files}).

\section{The four programs}
\label{sec:programs}

\index{programs}\index{window}\index{command line}\index{terminal, full screen}\index{web page}%
The same dictionary has four faces, which share the settings and the
history of searches:

\begin{longtable}{@{}lp{0.62\textwidth}@{}}
\toprule
Program & What it is\\
\midrule
\endhead
\texttt{words\_gui.py} & the window (section~\ref{sec:window})\\
\texttt{words.py} & the command line, with an interactive mode (section~\ref{sec:cli})\\
\texttt{words\_tui.py} & the window, but in a terminal, in full screen (section~\ref{sec:tui})\\
\texttt{words\_web.py} & a small web server, and the dictionary in the browser (section~\ref{sec:web})\\
\texttt{meanings.py} & \texttt{words.py} with only the dictionary forms and the meanings\\
\texttt{words\_edit\_gui.py} & the dictionary editor, in a window (section~\ref{sec:editing})\\
\texttt{words\_edit.py} & the dictionary editor, in the terminal\\
\texttt{words\_check.py} & checks the data files after they are edited by hand\\
\texttt{words\_install.py} & adds the program to the menu of the system\\
\bottomrule
\end{longtable}

\noindent
From the folder of the program they are run with \texttt{python3}:

\begin{verbatim}
$ python3 words_gui.py                     (a window)
$ python3 words_gui.py perspicio           (looking it up)
$ python3 words_gui.py -e look             (a window, to Latin)
$ python3 words.py rem acu tetigisti       (look up a line)
$ python3 words.py                         (interactive)
$ python3 words_tui.py                     (in the terminal)
$ python3 words_web.py                     (in the browser)
\end{verbatim}

\noindent\index{online version}%
The online version, \url{https://nivaca.pythonanywhere.com/}, is the page
of \texttt{words\_web.py} on a public server. Each visitor has their own
settings and earlier searches there, kept in the cookies of their browser,
and nothing of this is saved on the server.

\section{The window}
\label{sec:window}

\shot{gui_en}{1}{The window, in dark mode, with \textit{perspicio} looked up.}

\index{window}\index{search box}\index{look up}%
Figure~\ref{fig:gui_en} shows the window. At the top is the box where the
words are typed, and \menu{Look up}, which looks them up (\key{Return}
does the same). The box keeps the earlier searches: \key{Down}, or its
arrow, lists them, the most recent first, and choosing one looks it up
again.

\index{search in}\index{direction}\index{meanings in}\index{fuzzy search}%
Below it are the choices of the search:

\begin{description}
\item[Search in] \menu{Local}, this dictionary, or \menu{Logeion}, the
  dictionaries of the University of Chicago (section~\ref{sec:logeion}),
  with the list of Logeion's dictionaries beside it;
\item[the direction] \menu{Latin\,$\rightarrow$\,English} or
  \menu{English\,$\rightarrow$\,Latin} (with the Spanish meanings,
  \menu{Spanish} instead of \menu{English});
\item[Meanings in] the language of the meanings
  (section~\ref{sec:languages});
\item[Fuzzy search] when a word finds nothing, the results of its near
  spellings (section~\ref{sec:fuzzy}).
\end{description}

\noindent
These choices are kept: the next time, any of the programs starts as it
was left.

\index{results!links}\index{results!explanations}\index{double click}\index{right click}%
The results take the rest of the window. They are the same text as on the
command line, laid out for the eye: the word looked up is bigger, the
dictionary form is in bold with its headword in colour, the codes and the
notes are in grey, and the analyses of the forms in blue. The first word
of an entry is a link that looks it up, and a double click on any word
looks it up too (a word of the meanings is looked up in Latin, in the
other direction). With the mouse over the codes in brackets or over an
analysis (\texttt{GEN.PL}), a small box explains what they mean. The
right click copies or looks up a word.

\index{status line}\index{dark mode}%
The line at the bottom says what the program is doing. At its right, the
moon makes the windows dark (the dictionary, its dialogues and the
editor), and then the sun makes them light again. Beside it, the code of
the language of the interface (EN, ES or LA) opens the list of the
languages.

\subsection{The menus}

\index{menus}%
\begin{description}
\item[Dictionary] \menu{Edit dictionary\ldots} opens the editor
  (section~\ref{sec:editing}); \menu{Settings\ldots} shows all the
  settings (section~\ref{sec:settings}); \menu{Quit}.
\item[Edit] \menu{Cut}, \menu{Copy} and \menu{Paste}, in the box of the
  search or in the results.
\item[View] \menu{Find in results\ldots}, the size of the text, the
  fonts, and \menu{Dark mode}.
\item[Help] \menu{Help}, the guide to searching and to the results;
  \menu{Web page}, the page of the program; \menu{About}, the version and
  the credits.
\end{description}

\subsection{Finding a text in the results}

\index{find in results}%
\menu{View\then Find in results\ldots} (\key{Ctrl+F}) opens a small bar
under the results. What is typed there is highlighted as it is typed;
\key{Return} goes to the next place, \key{Shift+Return} to the previous
one, and \key{Esc} closes the bar.

\subsection{Text and fonts}

\index{fonts}\index{text size}\index{hi-res displays}%
\key{Ctrl++} and \key{Ctrl+-} make the text bigger and smaller, and
\key{Ctrl+0} brings back the usual size. \menu{View\then Font\ldots}
chooses the fonts: a fixed-width one for the results, because their
columns must stay aligned, and another for the rest of the window. Both
are kept. On a hi-res display the windows follow the scaling of the
desktop.

\section{Searching}
\label{sec:searching}

\subsection{From Latin}

\index{searching!from Latin}\index{capitals}\index{u and v, i and j}\index{long marks}%
Type a Latin word, or a whole line of Latin: each word is analysed in
turn. Capitals do not matter, nor do u and v, or i and j (\textit{iulius},
\textit{Julius} and \textit{IVLIVS} are the same word), nor long marks and
other accents: \textit{amīcus} is \textit{amicus}, and \textit{cælum} is
\textit{caelum}.

\index{tricks}\index{prefixes and suffixes}\index{enclitics}\index{medieval spellings}%
A word is first looked for as a stem and an ending of the dictionary.
When this finds nothing, the program tries prefixes, suffixes and
enclitics (\textit{-que}, \textit{-ne}, \textit{-ve}), and then other
spellings: a syncopated perfect, two words written as one
(\textit{rempublicam}), a final \textit{-st} for \textit{est}
(\textit{factumst}), and the spellings of medieval Latin (\textit{e} for
\textit{ae}, \textit{y} for \textit{i}, \textit{michi} for \textit{mihi}).
It always says what it did:

\begin{verbatim}
hystorie
  spelling: Medieval hystorie -> historiae
    'hystorie' may be a medieval spelling of historiae
  historia, historiae  noun, 1st declension, feminine  [XXXAX]
    a narrative of past events, history; ...
    histori-ae     GEN.SG
\end{verbatim}

\noindent\index{names}%
An unknown word with a capital is taken as a name, and left alone.

\subsection{To Latin}

\index{searching!to Latin}\index{parts of speech}%
In the other direction (\menu{English\,$\rightarrow$\,Latin} in the
window, \texttt{-e} on the command line), type one word of the language of
the meanings: the Latin entries whose meanings have it are listed, the
best first, six at most, and then how many more there are. A part of
speech after the word keeps only that one: \texttt{friend ADJ}. The parts
of speech are \texttt{N}, \texttt{PRON}, \texttt{ADJ}, \texttt{NUM},
\texttt{ADV}, \texttt{V}, \texttt{PREP}, \texttt{CONJ} and
\texttt{INTERJ}. Very common words (\textit{the}, \textit{a}) are never
looked for, and in Spanish the accents do not matter: \textit{corazon}
finds \textit{corazón}.

\subsection{The fuzzy search}
\label{sec:fuzzy}

\index{fuzzy search}\index{typos}%
Usually a word with a typo finds nothing. With the fuzzy search, a word of
three letters or more that still finds nothing is followed by the results
of the words one change away from it (a letter left out, added or
changed, or two letters swapped), three of them at most, those of the
most frequent entries first:

\begin{verbatim}
amicuss
  (unknown)
  (near spellings: amicus)

amicus
  amicus, amici  noun, 2nd declension, masculine  [XXXBO]
  ...
\end{verbatim}

\noindent
The search to Latin does the same with the words of the meanings:
\textit{freind} finds \textit{friend}, and \textit{wolfs} finds
\textit{wolf}. It is on unless it is turned off: \menu{Fuzzy search} in
the window and on the web page, \key{F3} in the terminal, or
\texttt{--no-fuzzy} on the command line, for one run.

\section{Reading the results}
\label{sec:results}

\index{results}%
This is what \texttt{py-words} shows for \textit{amicorum}:

\begin{verbatim}
amicorum
  amicus, amici  noun, 2nd declension, masculine  [XXXBO]
    A friend, patron, protector, companion; ... (Lewis and Short)
    amīc-ōrum      GEN.PL
  amicus, amica -um, amicior -or -us, amicissimus -a -um
      adjective, 1st/2nd declension  [XXXAO]
    friendly, kind, amicable, favorable, ... (Lewis and Short)
    amīc-ōrum      GEN.PL.M
    amīc-ōrum      GEN.PL.N
\end{verbatim}

\index{results!lines of}\index{principal parts}\index{long vowels}%
\begin{itemize}
\item The first line is the word looked up. Each entry it can be follows,
  indented.
\item The first line of an entry has its dictionary form (the principal
  parts: \textit{amicus, amici}), its part of speech and grammar in words
  (noun, 2nd declension, masculine), and the codes in brackets
  (section~\ref{sec:codes}); some entries add their age or frequency in
  words (archaic, lesser\ldots).
\item The next line is the meaning. Its senses are separated by
  semicolons, the commonest first. At its end, the dictionary it comes
  from, when it is not WORDS: (Lewis and Short), (De Miguel),
  (Echave-Sustaeta). \texttt{[en]} before a meaning says that it has no
  translation, and is given in English.
\item Then come the forms the word can be: the stem and the ending,
  joined by a hyphen, with the long vowels marked where they are known
  (\textit{amīc-ōrum}), and what the form is, in the abbreviations of
  section~\ref{sec:glosses}. One spelling can be several forms: the
  \textit{amīc-ōrum} of the adjective is masculine and neuter.
\item A line before an entry may say how the word was found: a prefix, a
  suffix or an enclitic, with its meaning, or a change of spelling.
\end{itemize}

\noindent\index{compound tenses}%
A participle followed by a form of \textit{esse} is read as a compound
tense:

\begin{verbatim}
laudatus
  laudo, laudare, laudavi, laudatus  verb, 1st conjugation  [XXXAX]
    to praise, laud, commend, extol, ... (Lewis and Short)
    laudāt-us      PRF.PASS.PTCP.NOM.SG.M
    participle + est  3SG.PRF.PASS.IND
      PERF PASSIVE PPL + verb TO_BE => PASSIVE perfect system
\end{verbatim}

\noindent\index{results!notes}%
Some notes go in parentheses:

\begin{longtable}{@{}>{\raggedright}p{0.4\textwidth}p{0.54\textwidth}@{}}
\toprule
Note & What it says\\
\midrule
\endhead
(unknown) & the word was not found\\
(near spellings: \ldots) & the results of the fuzzy search follow\\
(less likely forms are not shown) & the setting \texttt{trim\_output} left some out\\
(no match) & the search to Latin found nothing\\
(near words: \ldots) & the fuzzy search to Latin\\
(\ldots\ more not shown) & after the best six entries, to Latin\\
\bottomrule
\end{longtable}

\subsection{The codes in brackets}
\label{sec:codes}

\index{codes in brackets}\index{age}\index{area}\index{region}\index{frequency}\index{sources}%
The five letters in brackets after an entry say when, where, how often and
on whose authority the word is used, in this order: its age, its area (the
subject it belongs to), its region, its frequency and its source. An X
says only that there is nothing particular to say. So the
\texttt{[XXXBO]} of \textit{amicus} is read: in use throughout, general,
everywhere, frequent, and taken from the Oxford Latin Dictionary. An entry
can have more than one source: the others follow the five letters, each
after a \texttt{+} (\texttt{[XXXBO+A]} is also from Echave-Sustaeta). The
age, when it is known, and the frequency, when it is lesser or rarer, are
also written in words after the codes. The letters are these:

\begin{longtable}{@{}lp{0.8\textwidth}@{}}
\toprule
Letter & Age\\
\midrule
\endhead
X & in use throughout\\
A & archaic\\
B & early (before the classical period)\\
C & classical\\
D & late (after 200)\\
E & later (after 600)\\
F & medieval (after 1000)\\
G & scholarly (Renaissance and after)\\
H & modern\\
\bottomrule
\end{longtable}

\begin{longtable}{@{}lp{0.8\textwidth}@{}}
\toprule
Letter & Area\\
\midrule
\endhead
X & general\\
A & agriculture, plants, animals, the land\\
B & biology, medicine, the body\\
D & drama, music, the arts\\
E & the Church, the Bible\\
G & grammar, rhetoric, logic, literature, school\\
L & law, government, money, titles\\
P & poetry\\
S & science, philosophy, mathematics, measures\\
T & technical, architecture, surveying\\
W & war, the army, the navy\\
Y & mythology\\
\bottomrule
\end{longtable}

\begin{longtable}{@{}lp{0.8\textwidth}@{}}
\toprule
Letter & Region\\
\midrule
\endhead
X & everywhere\\
A & Africa\\
B & Britain\\
C & China\\
D & Scandinavia\\
E & Egypt\\
F & France, Gaul\\
G & Germany\\
H & Greece\\
I & Italy, Rome\\
J & India\\
K & the Balkans\\
N & the Netherlands\\
P & Persia\\
Q & the Near East\\
R & Russia\\
S & Spain\\
U & Eastern Europe\\
\bottomrule
\end{longtable}

\begin{longtable}{@{}lp{0.8\textwidth}@{}}
\toprule
Letter & Frequency\\
\midrule
\endhead
X & unknown\\
A & very frequent\\
B & frequent\\
C & common\\
D & lesser\\
E & uncommon\\
F & very rare\\
I & only in inscriptions\\
M & only in graffiti\\
N & only in Pliny\\
\bottomrule
\end{longtable}

\begin{longtable}{@{}lp{0.8\textwidth}@{}}
\toprule
Letter & Source\\
\midrule
\endhead
X & general or unknown\\
A & Echave-Sustaeta, \textit{Lengua latina. Vocabulario básico}\\
B & Beeson, \textit{Medieval Latin}\\
C & \textit{Cassell's Latin Dictionary}\\
D & Adams, \textit{Latin Sexual Vocabulary}\\
E & Stelten, \textit{Ecclesiastical Latin}\\
F & Deferrari, St.~Thomas Aquinas\\
G & Gildersleeve and Lodge, \textit{Latin Grammar}\\
H & Collatinus (Ouvrard)\\
I & Leverett, \textit{Lexicon}\\
J & Bracton, \textit{De Legibus}\\
K & \textit{Calepinus Novus} (Licoppe)\\
L & Lewis, \textit{Elementary Latin Dictionary}\\
M & Latham, \textit{Medieval Word List}; DMLBS\\
N & Nelson, word list\\
O & \textit{Oxford Latin Dictionary}\\
P & Souter, \textit{Later Latin}\\
Q & other dictionaries\\
R & Plater and White, \textit{Vulgate grammar}\\
S & Lewis and Short, \textit{A Latin Dictionary}\\
T & found in a translation\\
U & Du Cange\\
V & Blatt, Saxo\\
W & Whitaker's guess\\
Y & temporary\\
Z & sent by a user\\
m & De Miguel, \textit{Nuevo diccionario latino-español etimológico}\\
\bottomrule
\end{longtable}

\subsection{The abbreviations of the analyses}
\label{sec:glosses}

\index{glossing abbreviations}\index{Leipzig Glossing Rules}%
The forms are analysed with the abbreviations of the Leipzig Glossing
Rules, joined by full stops; the person and the number go together, so
\texttt{2SG.PRF.ACT.IND} is the 2nd person singular perfect active
indicative. Those marked * are not in the Leipzig list. With the setting
\texttt{use\_glosses} off (\texttt{--no-glosses} on the command line), the
analyses are written in words instead: \enquote{accusative singular}.

\begin{longtable}{@{}lp{0.6\textwidth}@{}}
\toprule
Abbreviation & Meaning\\
\midrule
\endhead
\texttt{NOM}, \texttt{GEN}, \texttt{DAT} & nominative, genitive, dative\\
\texttt{ACC}, \texttt{ABL}, \texttt{VOC}, \texttt{LOC} & accusative, ablative, vocative, locative\\
\texttt{SG}, \texttt{PL} & singular, plural\\
\texttt{M}, \texttt{F}, \texttt{N}, \texttt{C}* & masculine, feminine, neuter, common\\
\texttt{1}, \texttt{2}, \texttt{3} & first, second, third person\\
\texttt{PRS}, \texttt{IMPF}*, \texttt{FUT} & present, imperfect, future\\
\texttt{PRF}, \texttt{PLUP}*, \texttt{FUT.PRF}* & perfect, pluperfect, future perfect\\
\texttt{ACT}*, \texttt{PASS} & active, passive\\
\texttt{IND}, \texttt{SBJV}, \texttt{IMP} & indicative, subjunctive, imperative\\
\texttt{INF}, \texttt{PTCP}, \texttt{SUP}* & infinitive, participle, supine\\
\texttt{POS}*, \texttt{CMPR}*, \texttt{SPRL}* & positive, comparative, superlative\\
\texttt{CARD}*, \texttt{ORD}* & cardinal, ordinal numeral\\
\texttt{DISTR}, \texttt{MULT}* & distributive, multiplicative numeral (numeral adverb)\\
\bottomrule
\end{longtable}

\section{The languages}
\label{sec:languages}

\subsection{The language of the meanings}

\index{meanings in}\index{Spanish meanings}\index{De Miguel}\index{Echave-Sustaeta}%
The meanings are in English or in Spanish, and so is the rest of the
results (\enquote{sustantivo, 2.ª declinación, masculino}); the glossing
abbreviations are the same in both. The Spanish meanings come from De
Miguel's dictionary wherever it has the word, as he wrote them in 1867,
with his spelling (\textit{razon}, \textit{á}), and they say so at their
end, \enquote{(De Miguel)}; a few come from Echave-Sustaeta; the others
were translated from the English. It is chosen with \menu{Meanings in} in
the window and on the web page, with \key{F5} in the terminal, or with
\texttt{--language es} on the command line, for one run. The search to
Latin goes from the language chosen.

\subsection{The language of the interface}

\index{language of the interface}%
The programs speak English, Spanish or Medieval Latin: their menus,
buttons, messages,
help and settings, the editors and the checker. It is chosen with the
button with its code (EN, ES or LA), at the bottom right of the window
and at the top of the web page, with \menu{Language of the interface} in
the settings of the window and of the web page, with \key{F6} in the
terminal (its code is at the right of the top line), or at the end of the \texttt{\#}
dialogue of the command line (\texttt{--interface es} or
\texttt{--interface la} for one run). The window and the web page change
at once. Changing the interface changes the meanings too, to the same
language, if the dictionary has it; there are no Latin meanings, so with
the Latin interface they stay as they were, and so do the notes of the
results, such as \enquote{(unknown)}; after that each can be chosen apart,
so the interface can be in Spanish and the meanings in English.

\section{Logeion}
\label{sec:logeion}

\index{Logeion}%
Instead of the entries of this dictionary, the programs can show those of
Logeion (\url{https://logeion.uchicago.edu}), the Latin and Greek dictionaries of
the University of Chicago, with the short definition, the frequency and
examples from its corpus. It needs the internet. It is chosen with
\menu{Search in} in the window and on the web page, \key{F8} in the
terminal, \texttt{\tl G} in the interactive command line, or
\texttt{--logeion} for one run.

\index{lemmas}%
Logeion does not analyse inflected forms. So the words are analysed here
first, and Logeion is asked for their lemmas: \textit{dixit} shows
\textit{dico}, and \textit{est} shows \textit{edo} and \textit{sum}. Each
word shows its analysis first, without the meanings, and then Logeion's
pages of its lemmas. While Logeion is chosen, the search goes from Latin
only, and the other direction, the language of the meanings, the fuzzy
search and the editor are greyed out.

\index{Logeion!dictionaries}\index{DMLBS}%
The entries can be those of one dictionary of Logeion, or of all of them;
until another is chosen, it is the DMLBS. It is chosen with the list
beside \menu{Search in}, with \key{F5} in the terminal (while Logeion is
chosen), with \texttt{\tl D} in the interactive command line, or with
\texttt{--dictionary} and its code:

\begin{longtable}{@{}lp{0.7\textwidth}@{}}
\toprule
Code & Dictionary\\
\midrule
\endhead
\texttt{all} & all the dictionaries\\
\texttt{ls} & Lewis \& Short\\
\texttt{lewis} & Lewis, \textit{Elementary Latin Dictionary}\\
\texttt{georges} & Georges\\
\texttt{gaffiot} & Gaffiot 2016\\
\texttt{dmlbs} & DMLBS (\textit{Dictionary of Medieval Latin from British Sources})\\
\texttt{lane} & LaNe\\
\texttt{ducange} & Du Cange\\
\texttt{lana} & Lana-Michelone\\
\texttt{frieze} & Frieze-Dennison (Vergil)\\
\texttt{sinicum} & Latino-Sinicum\\
\bottomrule
\end{longtable}

\noindent
Changing the dictionary shows the results again at once, without asking
Logeion again. In the window, the links of an entry to other words look
those words up, and the other links open in the browser.

\section{The command line}
\label{sec:cli}

\shot{cli_en}{1}{The command line, interactive, with \textit{perspicio}.}

\index{command line}%
\texttt{words.py} looks up the words given after it, as one line of Latin,
or, without words, starts the interactive mode (figure~\ref{fig:cli_en}).
\texttt{meanings.py} is the same as \texttt{words.py -m}.

\begin{verbatim}
$ python3 words.py rem acu tetigisti
$ python3 words.py -e friend --pos N
$ python3 words.py -f text.txt -o results.txt
$ python3 words.py --language es amicorum
$ python3 words.py --logeion arma
$ python3 words.py --dictionary gaffiot arma
\end{verbatim}

\index{command line!options}%
\begin{longtable}{@{}>{\raggedright}p{0.36\textwidth}p{0.58\textwidth}@{}}
\toprule
Option & What it does\\
\midrule
\endhead
\texttt{-f FILE} & looks up each line of FILE (\texttt{-} for the standard input)\\
\texttt{-o FILE} & writes the results to FILE\\
\texttt{-e} & from the language of the meanings to Latin\\
\texttt{--pos POS} & with \texttt{-e}, only this part of speech\\
\texttt{-m} & only the dictionary forms and the meanings\\
\texttt{--no-glosses} & the analyses in words, not abbreviations\\
\texttt{--no-fuzzy}, \texttt{--fuzzy} & without or with the fuzzy search\\
\texttt{--logeion}, \texttt{--local} & Logeion, or this dictionary\\
\texttt{--dictionary DICT} & Logeion's dictionary (section~\ref{sec:logeion})\\
\texttt{--dictionaries} & lists Logeion's dictionaries\\
\texttt{--language CODE} & the language of the meanings (\texttt{en}, \texttt{es})\\
\texttt{--languages} & lists the languages of the meanings\\
\texttt{--interface CODE} & the language of the interface\\
\texttt{--datadir DIR} & the data files in DIR\\
\texttt{--guide} & prints the help: how to read the results\\
\texttt{--config} & creates the configuration file and says where it is\\
\texttt{--clear-history} & forgets the earlier searches\\
\texttt{--version} & the version and the credits\\
\bottomrule
\end{longtable}

\noindent
The options that choose something (the glosses, the fuzzy search, the
source, the languages) do so for that run only; the settings keep what
they had.

\subsection{The interactive mode}

\index{command line!interactive}%
In the interactive mode, each line typed after the prompt \texttt{=>} is
looked up, unless it starts with one of these:

\begin{longtable}{@{}>{\raggedright}p{0.36\textwidth}p{0.58\textwidth}@{}}
\toprule
Line & What it does\\
\midrule
\endhead
\texttt{\tl E}, \texttt{\tl L} & English to Latin; Latin to English\\
\texttt{\tl G} & Logeion or this dictionary\\
\texttt{\tl D [DICT]} & Logeion's dictionary (alone: a list to choose from)\\
\texttt{@FILE} & looks up each line of FILE\\
\texttt{\#} & changes the parameters (section~\ref{sec:settings})\\
\texttt{!} & changes the developer modes\\
\texttt{\tl H} & the help: how to search and read the results\\
\texttt{\tl A} & about: the version and the credits\\
\texttt{?} & the list of these lines\\
\bottomrule
\end{longtable}

\noindent\index{history of searches}%
\key{Up} and \key{Down} bring back the earlier searches, also those of
earlier sessions and of the other programs. In the dialogues of
\texttt{\#} and \texttt{!}, \texttt{?} explains each setting. Two empty
lines, or \key{Ctrl+D}, quit. \key{Ctrl+C} does not quit: it stops what is
being done (a long list of results, a file), and the program waits for
the next line.

\subsection{The pager and the colours}

\index{pager}\index{colours}%
In a terminal, results longer than the screen open in the pager of the
system, \texttt{less} (or the one \texttt{\$PAGER} names), as with
\texttt{man}: \key{PgUp} and \key{PgDn} scroll them, \texttt{/} finds a
text, and \texttt{q} goes back to the prompt. The results are in colour,
in the colours of the terminal, as in the window. Results written to a
file or a pipe are plain, and \texttt{NO\_COLOR=1} turns the colours off.

\section{The terminal, in full screen}
\label{sec:tui}

\shot{tui_en}{1}{\texttt{words\_tui.py}, in a terminal.}

\index{terminal, full screen}%
\texttt{words\_tui.py} is the window in a terminal (figure~\ref{fig:tui_en}):
a line for searching at the top, the results below it, and the keys in
the last line. \key{Up} and \key{Down} bring back the earlier searches, and
\key{PgUp} and \key{PgDn} scroll the results. On Windows it needs the
package \texttt{windows-curses}.

\index{terminal, full screen!keys}%
\begin{longtable}{@{}>{\raggedright}p{0.3\textwidth}p{0.64\textwidth}@{}}
\toprule
Key & Action\\
\midrule
\endhead
\key{F1} & the help, and the credits\\
\key{F2}, \key{Ctrl+L} & Latin to English, or English to Latin\\
\key{F3}, \key{Ctrl+O} & the settings (\key{Space} changes one)\\
\key{F4}, \key{Ctrl+E} & the dictionary editor\\
\key{F5}, \key{Ctrl+G} & the language of the meanings (with Logeion, its dictionary)\\
\key{F6} & the language of the interface\\
\key{F7} & copies all the results\\
\key{F8} & this dictionary or Logeion\\
\key{F9} & clears the history\\
\key{F10}, \key{Ctrl+Q} & quits\\
\bottomrule
\end{longtable}

\section{The web page}
\label{sec:web}

\shot{web_en}{0.8}{\texttt{words\_web.py}, in the browser.}

\index{web page}%
\texttt{words\_web.py} starts a small web server and opens the browser at
\url{http://127.0.0.1:8000/}, or at another port if this one is busy
(figure~\ref{fig:web_en}). The page has the box of the search, the
direction, \menu{Meanings in}, \menu{Fuzzy search}, \menu{Search in}, and
the results as in the window, with the same links and explanations; below
them, the earlier searches, all the settings, the help and the credits.
The moon makes the page dark, and the sun light again; until it is first
used, the page follows the system. Just before the moon, a round button
with the code of the language of the interface (EN, ES or LA) opens the
list of the languages.

\index{web page!addresses}%
Each search has its own address, such as
\url{http://127.0.0.1:8000/?q=amicus+est} (\texttt{\&lang=en} for English
to Latin), so it can be kept as a bookmark. The server reads the
dictionary again when an editor changes it; the dictionary cannot be
edited from the browser.

\index{web page!options}\index{network}%
Only this computer can open the page, unless \texttt{--host} says
otherwise: with \texttt{--host 0.0.0.0} it can be used from a phone on the
same network. There are no passwords, so do this only on a network you
trust. \texttt{--port} chooses the port, \texttt{--no-browser} does not
open the browser, and \key{Ctrl+C} stops the server.

\section{Settings}
\label{sec:settings}

\index{settings}\index{parameters}%
\menu{Dictionary\then Settings\ldots} shows all the settings of the
window, with what each one does (the mouse over a setting shows its
explanation); a change is applied at once. They are the same as the
parameters (\texttt{\#}) and the developer modes (\texttt{!}) of the
command line, \key{F3} in the terminal, and the settings of the web page.
The dialogue also has \menu{Language of the interface}, \menu{Font\ldots},
\menu{Clear history} and \menu{Default settings}.

\index{parameters!what the results show}%
\begin{longtable}{@{}>{\raggedright}p{0.36\textwidth}p{0.58\textwidth}@{}}
\toprule
Parameter & What it does, when it is on\\
\midrule
\endhead
\multicolumn{2}{@{}l}{\itshape What the results show}\\
\texttt{show\_dictionary\_forms} & each entry with its principal parts and grammar\\
\texttt{show\_only\_meanings} & only the entries and their meanings, without the forms\\
\texttt{use\_glosses} & the analyses in glossing abbreviations, not in words\\
\texttt{show\_long\_vowels} & the long vowels marked in the forms (\textit{amīc-ōrum})\\
\texttt{show\_codes} & the codes in brackets (\texttt{[XXXBO]})\\
\texttt{explain\_codes} & the age and the frequency always in words, not only when unusual\\
\texttt{trim\_output} & leaves out the least likely results (and, to Latin, all but the best six)\\
\texttt{trim\_medieval} & with \texttt{trim\_output}, also medieval forms when others were found\\
\addlinespace
\multicolumn{2}{@{}l}{\itshape How words are looked up}\\
\texttt{fuzzy\_search} & the near spellings of a word that finds nothing\\
\texttt{use\_affixes} & an unknown word read as a known one with a prefix or suffix\\
\texttt{use\_tricks} & other spellings, and two words written as one\\
\texttt{medieval\_spellings} & with \texttt{use\_tricks}, the medieval spellings too\\
\texttt{read\_compounds} & a participle and a form of \textit{esse} as a compound tense\\
\texttt{ignore\_unknown\_names} & no tricks on an unknown word with a capital\\
\texttt{near\_entries} & for an unknown word, the entries near where it would be\\
\addlinespace
\multicolumn{2}{@{}l}{\itshape Developer modes}\\
\texttt{show\_entry\_source} & which dictionary each entry is from, and its number there\\
\texttt{pearse\_codes} & each line of the results starts with a code (for other programs)\\
\texttt{pause\_in\_screen\_output} & long results in the pager, on the command line\\
\bottomrule
\end{longtable}

\noindent\index{trim\_output}%
Poets use unusual words and forms, so \texttt{trim\_output} can leave out
the right one; with a verse that gives strange results, turn it off.

\section{Editing the dictionary}
\label{sec:editing}

\index{dictionary editors}%
Two editors add, correct and delete the entries of the dictionary,
\texttt{data/dictionary.xml}: one in a window, \texttt{words\_edit\_gui.py},
which \menu{Dictionary\then Edit dictionary\ldots} also opens, and one in
the terminal, \texttt{words\_edit.py}. In both of them:

\begin{itemize}
\item an entry is found by a Latin form (\textit{amicorum}), a stem or a
  word of its meaning;
\item it is corrected, or deleted, or a new one is made, empty or like
  the one shown (\menu{New like this}); the fields depend on the part of
  speech, and the codes are chosen from lists that explain them;
\item the preview shows the entry as the results will show it, and any
  problem it has; \menu{Try a word} shows what the program makes of a word
  with the dictionary as it is now;
\item the meanings in Spanish can be edited too, beside the English ones,
  after choosing the language with \menu{Meanings in} (\key{F7} in the
  terminal).
\end{itemize}

\index{dictionary editors!saving}\index{backup files}%
A change goes into the dictionary with \menu{Apply changes} (\key{F3} in
the terminal), and the file is written only with \menu{Save} (\key{F2}).
Leaving an entry with changes not applied asks what to do with them. The
first save of a session keeps the original files as \texttt{.bak} files.

\index{dictionary editors!keys}%
\begin{longtable}{@{}>{\raggedright}p{0.3\textwidth}p{0.64\textwidth}@{}}
\toprule
Key (terminal) & Action\\
\midrule
\endhead
\key{F3}, \key{Ctrl+A} & applies the changes (or adds the new entry)\\
\key{F4}, \key{Ctrl+N} & a new entry\\
\key{F5}, \key{Ctrl+D} & a new entry like this one\\
\key{F6}, \key{Ctrl+T} & tries a word\\
\key{F7}, \key{Ctrl+L} & the language of the meanings edited besides English\\
\key{F8}, \key{Ctrl+K} & deletes the entry\\
\key{F2}, \key{Ctrl+S} & saves the file (applying the changes first)\\
\key{F10}, \key{Ctrl+Q} & quits\\
\key{Tab} & the next area: the search, the list, the form\\
\key{Ctrl+F} & find\\
\bottomrule
\end{longtable}

\noindent\index{dictionary editors!read-only data}\index{data folder}%
In the AppImage and the RPM the data cannot be written, so the editors are
off. To edit the dictionary there, copy the folder \texttt{data} somewhere,
and give the copy as the data folder: in the configuration file
(\texttt{folder} in \texttt{[data]}), with \texttt{--datadir}, or with
\texttt{\$PY\_WORDS\_DATADIR}.

\index{words\_check.py}%
The files can also be edited by hand, in any text editor: copy the entry
of a similar word and change it. Then \texttt{python3 words\_check.py}
reports, with the line numbers, what the program cannot read and the
entries that probably do not do what was meant.

\section{Keys of the window}
\label{sec:keys}

\index{keys}%
\begin{longtable}{@{}>{\raggedright}p{0.4\textwidth}p{0.54\textwidth}@{}}
\toprule
Key & Action\\
\midrule
\endhead
\key{Return} & look up\\
\key{Down} (in the box of the search) & the earlier searches\\
\key{Ctrl+F} & find in the results\\
\key{Return}, \key{Shift+Return} (finding) & the next and the previous place\\
\key{Esc} (finding) & close the bar of the search\\
\key{Ctrl++}, \key{Ctrl+-} & larger and smaller text\\
\key{Ctrl+0} & the usual size\\
\key{Ctrl+X}, \key{Ctrl+C}, \key{Ctrl+V} & cut, copy, paste\\
\key{F1} & help\\
\key{Ctrl+Q} & quit\\
\key{Esc} (in a dialogue) & close it\\
\bottomrule
\end{longtable}

\section{Where things are kept}
\label{sec:files}

\index{configuration files}\index{config.toml}\index{history of searches}\index{cache}%
The settings, the fonts, the dark mode and the choices of the window are
kept in one configuration file, \texttt{config.toml}. It is plain text,
with a comment for each setting, and it can be edited by hand;
\texttt{python3 words.py --config} creates it and says where it is. The
program also keeps the last 1,000 searches, and a cache of the dictionary
that makes it start in a quarter of a second. All of them are in the
user's folders:

\begin{longtable}{@{}llp{0.6\textwidth}@{}}
\toprule
What & System & Where\\
\midrule
\endhead
Settings & Linux & \texttt{\textasciitilde/.config/py-words/config.toml}\\
 & macOS & \texttt{\textasciitilde/Library/Application Support/py-words}\\
 & Windows & \texttt{\%APPDATA\%\textbackslash py-words}\\
\addlinespace
History & Linux & \texttt{\textasciitilde/.local/state/py-words/history}\\
 & macOS & \texttt{\textasciitilde/Library/Application Support/py-words}\\
 & Windows & \texttt{\%LOCALAPPDATA\%\textbackslash py-words}\\
\addlinespace
Cache & Linux & \texttt{\textasciitilde/.cache/py-words}\\
 & macOS & \texttt{\textasciitilde/Library/Caches/py-words}\\
 & Windows & \texttt{\%LOCALAPPDATA\%\textbackslash py-words}\\
\bottomrule
\end{longtable}

\noindent\index{environment variables}%
\texttt{\$PY\_WORDS\_CONFIG}, \texttt{\$PY\_WORDS\_HISTORY} and
\texttt{\$PY\_WORDS\_CACHE} give other places; an empty
\texttt{\$PY\_WORDS\_HISTORY} keeps no history, and an empty
\texttt{\$PY\_WORDS\_CACHE} no cache. The cache can always be deleted.
The data files are looked for first with \texttt{--datadir}, then in
\texttt{\$PY\_WORDS\_DATADIR}, then in the folder the configuration file
names, and last in the folder \texttt{data} next to the program.

\section{When something goes wrong}

\index{troubleshooting}%
\begin{description}
\item[The program says that Python is too old] install Python 3.13 or
  later, or use the AppImage.
\item[The window does not open: no module named tkinter] install Tk
  (section~\ref{sec:install}).
\item[A word is not found] check \menu{Fuzzy search}, and the settings
  \texttt{use\_tricks} and \texttt{medieval\_spellings}; an unknown word
  with a capital is taken as a name, unless
  \texttt{ignore\_unknown\_names} is off.
\item[A verse gives strange results] turn off \texttt{trim\_output}, which
  leaves out unusual forms.
\item[Logeion does not answer] it needs the internet; \menu{Search in\then
  Local} goes back to this dictionary.
\item[\menu{Edit dictionary\ldots} is greyed out] the data cannot be
  written (AppImage, RPM), or Logeion is chosen
  (section~\ref{sec:editing}).
\item[The text of the window is too small or too big] \key{Ctrl++} and
  \key{Ctrl+-}, or \menu{View\then Font\ldots}.
\item[A setting is wrong] delete \texttt{config.toml}
  (section~\ref{sec:files}); the program starts with its defaults.
\end{description}

\noindent
Problems and wishes can be written as issues on GitHub or Codeberg (the
addresses are on the first page), or sent to
\href{mailto:nicolasvaughan@proton.me}{\texttt{nicolasvaughan@proton.me}}.

\printindex

\end{document}

% ===============================================
%%% Local Variables:
%%% mode: latex
%%% lsp-ltex-plus-language: "en-GB"
%%% TeX-master: t
%%% End:
